% Registro de lo recortado de cap3/3_2_validacion.tex (no se carga con \input).

% ===================================================================
% RECORTADO de cap3/3_2_validacion.tex: §4.2.3 Compatibilidad con registros anteriores
% ===================================================================
\subsection{Compatibilidad con registros anteriores}
\label{sec:val_compatibilidad}

El formato del registro de sesión descrito en la
\cref{sec:formato_datos} no fue el primero. Las sesiones anteriores a la
integración del módulo de temperatura se guardaron con un conjunto
menor de columnas, y algunas de ellas conservan valor documental.

La lectura de sesiones se implementó tratando el conjunto de campos
declarado en \texttt{COLUMNAS\_REQUERIDAS} como un subconjunto del
encabezado presente en el archivo, en lugar de exigir coincidencia
exacta. Con ese criterio, un registro de dieciocho columnas se abre sin
modificación y sin pérdida de los campos que sí contiene. La
verificación se realizó sobre sesiones reales conservadas del formato
anterior.

El criterio no es una previsión abstracta: el encabezado cambió cinco
veces a lo largo del proyecto. Las sesiones de mayo y junio de 2026
tienen 15 columnas (una de junio, 16); las de agosto, 18; las del 8 y
del 22 de septiembre, 21; las últimas del 22 de septiembre, 23, tras
sustituirse el campo de pulsos estimados por las columnas de
adquisición; y la versión congelada escribe 34, con las lecturas
individuales de los cuatro sensores y los límites de la ventana de
integración. Todas contienen el subconjunto declarado en
\texttt{COLUMNAS\_REQUERIDAS}.

La contraparte de esa decisión es que la ausencia de una columna deja de
ser un error de apertura y pasa a ser una condición que debe detectarse
más tarde, en la verificación del contenido. Ese desplazamiento es
deliberado y se examina en la subsección siguiente.


% ===================================================================
% RECORTADO de cap3/3_2_validacion.tex: §4.2.5 Propagación de advertencias a la interfaz
% ===================================================================
\subsection{Propagación de advertencias a la interfaz}
\label{sec:val_advertencias}

El sistema no interrumpe una secuencia ante una anomalía recuperable:
la registra y la señala. La razón es de economía experimental: un
barrido por temperatura sobre el enfriamiento pasivo de la muestra no
puede repetirse a voluntad, y abortarlo por una lectura perdida o un
objetivo rebasado descartaría los puntos que todavía pueden medirse;
la fila afectada queda marcada en el registro, como se ilustra en la
\cref{sec:falla_temperatura}. Esa decisión hace que el indicador de
advertencia de la interfaz sea el único aviso que recibe el operador
mientras la sesión transcurre.

La comprobación \texttt{probar\_advertencias\_gui.py} ejercita cuatro
trayectorias de esa señalización y las cuatro resultan conformes. El
alcance de esa cifra, sin embargo, es menor de lo que sugiere: la
comprobación calcula su código de retorno antes de ejecutar el único
paso que requiere confirmación visual del operador, de modo que ese
paso no puede hacerla reprobar. Los cuatro resultados corresponden a
las trayectorias automatizables; el paso restante queda como
comprobación asistida, y así se reporta. El resultado, además, se
obtuvo sobre la versión de la ventana principal anterior a la
auditoría de la \cref{sec:val_auditoria}, que reescribió buena parte
de ese módulo, y no consta que la comprobación se haya vuelto a correr
con la versión congelada.


% ===================================================================
% RECORTADO de cap3/3_2_validacion.tex: §4.2.6 Un defecto real encontrado al reconstruir las pruebas
% ===================================================================
\subsubsection*{Un defecto real encontrado al reconstruir las pruebas}

La reconstrucción no fue un trámite. Al correr las pruebas nuevas
contra la versión anterior a la auditoría, una de ellas quedó detenida
indefinidamente en un diálogo de selección de carpeta que el arnés no
interceptaba, y el arnés se corrigió. En la misma sesión de trabajo se
identificó un efecto de la propia auditoría que ningún commit había
documentado. Al
hacer que la interfaz siguiera en vivo la sesión activa del
almacenamiento, una sesión reabierta con «Abrir CSV…» pasaba a recibir
las capturas siguientes: un Guardar manual, o una secuencia
automática, agregaba filas a esa sesión ---posiblemente de otro día---
sin pasar por el diálogo de carpeta y sin aviso. La consecuencia es la
mezcla silenciosa de mediciones de días distintos en un mismo
registro.

La corrección (\texttt{f3a8c31}) pide una confirmación antes de la
primera escritura en una sesión reabierta, sea un Guardar o el inicio
de una secuencia. El diálogo nombra la sesión y su fecha de creación,
tomada de \texttt{metadatos\_sesion.txt} o, en su defecto, del prefijo
del identificador de sesión; si ninguna está disponible, lo declara en
lugar de inventarla. La respuesta por defecto es No, que cancela solo
esa escritura sin dejar fila ni archivo de muestras; un Sí vale para
el resto de esa reapertura, y reabrir cualquier registro o crear una
sesión nueva vuelve a exigir la confirmación. Una segunda corrección
(\texttt{c67ae06}) movió la pregunta para que aparezca después de
aceptar el inicio de la secuencia: en la primera versión, un Sí seguido
de la cancelación de la secuencia dejaba la reapertura confirmada y el
siguiente Guardar escribía sin preguntar.

La prueba correspondiente cuenta con doce casos y discrimina entre
diseños: contra un diseño que preguntara en cada Guardar fallan dos de
ellos, y el caso añadido con la segunda corrección falla contra la
primera. Además se realizó una corrida de punta a punta fuera del
repositorio, arrancando la aplicación como en uso normal, con clics
reales y un osciloscopio emulado en la capa \texttt{vxi11}: sin
reabrir ninguna sesión no aparece ninguna pregunta; tras «Abrir
CSV…», cancelar no escribe y aceptar agrega la fila, sin nuevas
preguntas en los Guardar siguientes; y reabrir el mismo registro
vuelve a preguntar. El inicio de una secuencia no se recorrió de punta
a punta, porque exige un láser conectado; lo cubre únicamente la
prueba automatizada.


% ===================================================================
% RECORTADO de cap3/3_2_validacion.tex: §4.2.7 Poblado de escalas; captura manual sin módulo de temperatura; limpieza de la bandera
% ===================================================================
\subsubsection*{Poblado de las columnas de escala}

La regresión descrita en la \cref{sec:falla_regresion} se corrigió
restituyendo las consultas \texttt{CH1:SCALE?} y
\texttt{HORizontal:SCAle?} dentro de la rutina de lectura del
preámbulo. El criterio fijado excluía explícitamente el valor
\num{0.0} como aprobación, pues es lo que escriben los manejadores de
excepción cuando la consulta se lanza y falla: solo un número
consistente con el panel del instrumento confirma la corrección. En la
sesión \texttt{20260908\_184623} las catorce capturas registran
\texttt{CH\_SCALE} $=\num{0.001}$ y \texttt{HOR\_SCALE} $=\num{1e-5}$,
equivalentes a \SI{1}{\milli\volt} y \SI{10}{\micro\second} por
división, en coincidencia con el panel. Conforme.

\subsubsection*{Captura manual con el módulo de temperatura
desconectado}

El criterio fijado exigía tres cosas de la fila resultante:
\texttt{error\_flag = 1}, un \texttt{error\_desc} que nombrara al
dispositivo ausente y un campo de temperatura en cero. La primera
ejecución cumplió únicamente la primera. Las capturas
\texttt{m0012}--\texttt{m0014} de la sesión \texttt{20260908\_184623},
tomadas en una secuencia por tiempo después de desconectar el módulo,
levantaron la bandera, pero escribieron \SI{24.25}{\celsius} sin sensor
conectado y una descripción genérica. El defecto y su corrección se
documentan en la \cref{sec:falla_fantasma}.

La corrección modificó además el propio criterio. El campo en cero se
sustituyó por \texttt{nan}, porque \SI{0}{\celsius} es una temperatura
físicamente válida y un análisis posterior no podría distinguirla de
una medición; \texttt{nan} no admite esa lectura. En la sesión
\texttt{20260908\_190614}, con el módulo desconectado, tanto la
captura manual \texttt{m0001} como las capturas de la secuencia que la
sigue registran \texttt{error\_flag = 1}, \texttt{temperatura = nan} y un
\texttt{error\_desc} que nombra al ESP32 y su puerto. Conforme bajo el
criterio corregido.

\subsubsection*{Limpieza automática de la bandera del monitor}

El criterio exigía que la bandera del monitor de conexión se limpiara
por sí sola en el siguiente ciclo de consulta tras restablecerse el
módulo. La primera ejecución tampoco lo cumplió: el registro de la
sesión \texttt{20260908\_190614} muestra que cada recuperación
(\texttt{m0005} y \texttt{m0013}) ocurre después de una pausa
prolongada, la que toma accionar el botón de reconexión en la
interfaz. No había limpieza automática porque no había reconexión
automática: el método correspondiente se limitaba a informar el estado
de la conexión sin intentar restablecerla.

Se implementó el reintento sobre el puerto configurado y, al recibirse
una trama válida, el reinicio del hilo de lectura. La sesión
\texttt{20260908\_200938} lo verifica: las capturas \texttt{m0003} y
\texttt{m0004} se registran con el módulo desconectado y
\texttt{temperatura = nan}, y \texttt{m0005} vuelve a
\texttt{error\_flag = 0} con lectura real, sin intervención del
operador.

La implementación reintenta únicamente sobre el puerto declarado en la
configuración y no recorre los puertos disponibles buscando el
dispositivo. La restricción es deliberada. Un puerto no se reasigna
solo: cuando cambia es porque alguien manipuló la conexión física o el
sistema operativo la renumeró tras esa manipulación, y en ambos casos
hay una persona presente que puede corregir la configuración. Recorrer
los puertos implicaría, en cambio, abrir puertos ajenos al módulo
---entre ellos el del láser--- para averiguar qué hay detrás de cada
uno, lo que introduce un riesgo mayor que el problema que resuelve.

El costo de esa decisión se materializó en la misma sesión: el módulo
apareció en \texttt{COM5} al inicio y en \texttt{COM4} al final, y esa
transición sí requirió editar la configuración y reiniciar la
aplicación. La reconexión automática cubre la pérdida de comunicación
sobre un puerto estable, no la reasignación del puerto. Ese costo se
redujo después: desde el 14 de septiembre, el diálogo «Conexión» de la
interfaz permite elegir el puerto del módulo, el del láser y la
dirección del osciloscopio, guarda la elección en un archivo de
configuración del usuario y la aplica en la siguiente conexión de cada
dispositivo, sin editar código ni reiniciar la aplicación.


% ===================================================================
% RECORTADO de cap3/3_2_validacion.tex: §4.2.7 Exportación a .mat; «Lo que estas cinco comprobaciones dejan ver»
% ===================================================================
\subsubsection*{Exportación al formato \texttt{.mat}}

La quinta comprobación se ejecutó sobre los datos de esa misma sesión.
El archivo resultante contiene el arreglo crudo, las muestras
convertidas a volts, el eje temporal reconstruido, la frecuencia de
muestreo, el número de puntos y los veintiún campos de metadatos del
registro de sesión. Conforme. El exportador copia todas las columnas
de la fila, de modo que con el formato congelado el archivo lleva las
34; esa versión del exportador no se ha ejercitado sobre una sesión
real, porque no existe todavía una sesión de 34 columnas.

\subsubsection*{Lo que estas cinco comprobaciones dejan ver}

Tres de las cinco resultaron conformes en su primera ejecución. Dos no,
y ninguna de las dos podía haberse detectado sin el instrumento: una
escribía un dato inexistente y la otra ofrecía una función que nunca
llegó a implementarse. Ambas corresponden a trayectorias que la
validación por simulación no alcanza por construcción y no por descuido
en su diseño: \texttt{generar\_sesion\_prueba.py} fabrica el archivo
que produce una desconexión, no la desconexión misma.

El resultado respalda la separación adoptada al inicio de esta sección
y marca al mismo tiempo su frontera. La simulación verificó
correctamente todo lo que estaba a su alcance; lo que quedó fuera de su
alcance falló en dos de cinco casos.


% ===================================================================
% RECORTADO de cap3/3_2_validacion.tex: «Familias de defectos corregidos» (pasada 2; sustituida por la tabla de clases)
% ===================================================================
\subsubsection*{Familias de defectos corregidos}

La auditoría se organizó en cuatro grupos, uno por subsistema: la
ventana principal (grupo~A), el osciloscopio (B), el módulo de
temperatura (C) y el almacenamiento (D), con un commit de corrección
por grupo. Los defectos encontrados se reparten en cinco familias.

\begin{enumerate}[leftmargin=*]
  \item \emph{Éxito anunciado por envío, no por confirmación.} El
    stop de emergencia y la aplicación de parámetros del láser
    informaban éxito sin consultar el resultado de cada comando; el
    estado de emisión se fijaba en marcha o detenido por haber enviado
    la orden; el canal del osciloscopio se daba por cambiado aunque el
    comando hubiera fallado. Tras la corrección, el estado del láser se
    marca como incierto hasta releer su registro de estado, y la
    selección de canal visible vuelve al canal real si el cambio
    falla.
  \item \emph{Señales de error sin destinatario.} Los avisos del
    módulo de temperatura (\texttt{TempWorker.\allowbreak error}), los fallos de
    guardado (\texttt{Almacenamiento.\allowbreak guardado\_err}) y los errores del
    hilo de medición (\texttt{MedicionWorker.\allowbreak error}) no estaban
    conectados a ningún elemento de la interfaz: un fallo al guardar
    no producía aviso visible. Se conectaron al registro de la
    ventana, igual que las del láser y el osciloscopio.
  \item \emph{Filas que parecen limpias sin serlo.} Un sensor
    DS18B20 ausente no activaba \texttt{error\_flag}, aunque el propio
    mensaje del módulo lo prometía; si la escritura del archivo
    \texttt{.npy} fallaba, la fila se registraba con
    \texttt{error\_flag = 0} y sin archivo de muestras; y la fila de
    una captura manual tomaba la temperatura y los parámetros del
    láser al pulsar Guardar, no al capturar. Las tres condiciones
    marcan ahora la fila o registran el valor del instante de la
    captura.
  \item \emph{Estado obsoleto.} Al agotarse los reintentos de lectura
    del osciloscopio, el controlador seguía apuntando al enlace caído;
    la estimación del tiempo de captura usaba la longitud de registro
    leída al conectar; el cierre de la ventana a mitad de secuencia
    dejaba el osciloscopio con el número de promedios de la secuencia;
    y el hilo de temperatura giraba indefinidamente ante una racha de
    tramas inválidas, porque solo contaba como silencio la ausencia
    total de líneas.
  \item \emph{Liberación prematura de la interfaz.} El stop de
    emergencia liberaba los controles antes de que la secuencia
    confirmara su fin, lo que permitía arrancar una secuencia nueva con
    la anterior todavía activa.
\end{enumerate}

